\section{Evaluation Scope and Evidence Map}
\label{sec:supp-evaluation-scope}

\begin{figure}[t]
\centering
\includegraphics[width=\columnwidth,height=0.60\textheight,keepaspectratio]{figures/fig_supplement_first_page_evidence_quickstart.png}
\caption{First-page evidence quickstart. The card opens with proxy,
GRScenes target-view, material-effect, navigation media, and opening-render
examples before the detailed appendix pages. The compass panel summarizes the
reading order used throughout the supplement.}
\label{fig:supp-first-page-evidence-quickstart}
\end{figure}

The main paper uses four evaluation checks: proxy preservation, GRScenes VLM
grounding, material-effect diagnostics, and embodied-stack usability. The
purpose of this appendix is to make the scope explicit. Each check answers a
different question. Keeping them separate makes clear which measurement surface
supports each result.

\begin{figure*}[p]
\centering
\includegraphics[width=\textwidth,height=0.82\textheight,keepaspectratio]{figures/fig_supplement_visual_evidence_roadmap.png}
\caption{Evidence roadmap. The lanes front-load representative render
panels from the proxy, GRScenes VLM, material-effect, and InternNav sections so
that the rest of the supplement can be read as an evidence map rather than as a
sequence of disconnected tables. Each lane points readers to the detailed
evidence pages that follow.}
\label{fig:supp-visual-evidence-roadmap}
\end{figure*}
\clearpage

\begin{table*}[t]
\centering
\footnotesize
\setlength{\tabcolsep}{3pt}
\begin{tabular}{p{0.15\textwidth}p{0.30\textwidth}p{0.24\textwidth}p{0.22\textwidth}}
\toprule
Gate & Evidence used & Claim supported & Forbidden promotion \\
\midrule
Proxy similarity & Four Isaac Sim assets; 24 matched render pairs; PSNR, SSIM, LPIPS, CLIP, DINOv2, and load/render diagnostics. & Converted assets pass a visual and feature-screening gate. & Task preservation, broad robustness, or speedup. \\
VLM grounding & GRScenes 15-pair clean pilot and frozen 30-pair target-centered stress set with answer, point-in-box, and coordinate scoring. & Material perturbation can be evaluated as bounded answer and point-grounding evidence. & Population-level GRScenes benchmark claims or general VLM rankings. \\
Material mechanisms & Selected opacity/transparency, emission, normal/bump, and displacement/height bins, plus clearcoat and procedural supplemental cases; original/noMDL/NVIDIA conditions are separated. & Covered bins have selected static and qualitative support; clearcoat and procedural texture remain scoped failure or limitation cases. & NVIDIA population failure rates, official-scene NVIDIA performance, or procedural-texture success. \\
Embodied-data sanity & Official KuJiaLe InternNav 99 paired episodes over three scenes plus official-scene load/render stability. & Converted scenes can enter the selected official downstream stack and load/render repeatedly. & Broad embodied-navigation robustness, all-scene preservation, or noMDL speedup. \\
\bottomrule
\end{tabular}
\caption{Evidence-gate registry for the ACL/ARR claim boundary. Each gate lists the evidence that may be used, the claim it supports, and the promotion that remains forbidden without new evidence.}
\label{tab:acl_evidence_gate_registry}
\end{table*}

\clearpage

Table~\ref{tab:acl_evidence_gate_registry} gives the shortest reading guide.
Proxy render scores cover matched cameras and lighting. VLM rows cover the
specified prompt, coordinate format, and target set. Material panels make
selected failure modes visible. The InternNav run checks whether matched scenes
can be exercised by the official stack. Presenting the same result as a table,
derivation, and visual example is useful here because the paper uses several
measurement surfaces: metrics, prompts, material cases, and navigation media.

\subsection{Reading Order}

Appendix~\ref{sec:supp-render-atlas} starts with the render atlas
so readers can inspect the actual image evidence before the formal details.
Appendix~\ref{sec:supp-derivations} gives the mathematical definitions used by
the evaluation checks. Appendix~\ref{sec:supp-vlm-protocol} expands the prompt and
coordinate protocol. Appendix~\ref{sec:supp-grscenes-visuals} collects
additional GRScenes diagnostics. Appendix~\ref{sec:supp-material-effects}
summarizes material-effect and NVIDIA-baseline risk cases.
Appendix~\ref{sec:supp-internnav-visuals} is the largest navigation section; it
shows selected stills and explains the media inventory.
Appendix~\ref{sec:supp-theory} records mechanism-level hypotheses and
failure-mode interpretations. Appendix~\ref{sec:supp-reproducibility} lists the
artifact and media sources. The appendix supports the main-paper results;
broader claims about lossless MDL semantics, new VLM backends,
population-level material rates, or general navigation robustness would require
separate experiments.
